\section{Related Work}
\subsection{Comparable Work Is Only One Part of a Comparable Measurement}
Benchmarking frameworks establish what a system is asked to execute and how performance and energy should be reported. MLPerf Tiny and MLPerf Power address widely different operating scales, while MLPerf Inference specifies power-measurement procedures~\cite{banbury2021mlperftiny,tschand2024mlperfpower,mlcommonsInferencePower}. Edge studies apply such measurements to specialised platforms and deployment choices~\cite{liang2020,baller2021,libutti2020usbaccelerators}. These efforts make experimental work more comparable. They do not make a module-input observation, a card-input observation and a supply-input observation interchangeable. Our concern is the measurement decision that precedes interpreting their numbers: which quantity and interval does a particular acquisition actually support?

\subsection{Calibration and Temporal Resolution Answer Different Questions}
External instruments provide a way to test internal readings. Shalavi et al. use externally referenced regression to calibrate built-in sensors on several Jetson devices~\cite{shalavi2023jetsoncalibration}. PowerSensor3 provides open measurement hardware and software for direct power acquisition~\cite{vanderVlugt2025powersensor3}. These approaches address the reliability and accessibility of the observation. A further test is needed when the desired output changes from a long average to a short event: agreement in level does not establish sufficient temporal detail.

The closest methodological connection is Gralka et al.'s \emph{Power Overwhelming}~\cite{gralka2024poweroverwhelming}. It combines oscilloscope measurements, external instruments and vendor interfaces, with synchronisation to rendering activity. Its comparison of aggregate readings and within-frame structure demonstrates why one instrument can be useful for energy totals yet inadequate for detailed activity analysis. We build on that distinction rather than treating the combination of meters as new. Our focus is the joint interpretation of finite-interval reconstruction, recorded spectral coverage and practical measurement boundaries in edge-AI inference.

\subsection{Telemetry Polling Is Not Physical Sampling}
Software-accessible power data introduce update and averaging behaviour in addition to electrical response. Evaluations of Jetson sensors, RAPL counters and software-based meters show why the interface and reported domain belong in the measurement description~\cite{aslan2022jetsonbuiltin,alt2024raplheteromemory,jay2023softwarepowermeters}. Yang et al. examine NVIDIA sensor behaviour across GPU generations, motivating device-specific rather than universal assumptions about built-in telemetry~\cite{yang2024parttime}.

Dauner et al. study how the polling interval affects RAPL and NVML energy readings, including comparisons with external instruments~\cite{dauner2026frequency}. Their results concern the interaction of counters, stale readings and acquisition overhead. Our same-trace experiment changes a numerical grid after the physical execution has finished; our deployed-meter experiment instead preserves each path's recorded intervals and response. Keeping these experiments separate avoids explaining every difference between physical acquisitions as a sampling-grid effect.

Together, this literature motivates a hierarchy of checks, not a single winning meter. Calibration addresses the level, reconstruction tests a defined integral, spectral analysis describes observed dynamics, and physical repetitions test the deployed procedure. The present study connects these checks through Jetson and Hailo observations. Its duration-wise pairing adds a practical distinction between agreement in mean power and agreement in integrated energy without claiming a synchronised, event-level telemetry calibration.
